\section{Detection of Subfunctionalization and Dosage Effects via Dynamic Programming}

    This section formalizes the algorithmic strategies to identify subfunctionalization and dosage effect patterns among in-paralogous gene sets, based on semantic expression vectors in Tensor Omics. The procedures leverage dynamic programming and bitmask representations to achieve efficient, branch-pruned exploration of possible paralog subsets.

    \subsection{Overview}

        Given:
        \begin{itemize}
            \item The ancestral ortholog expression vector $\vec{o}$.
            \item A set of $n$ in-paralog expression vectors $\vec{p}_1, \vec{p}_2, \dots, \vec{p}_n$, stored as columns of a Fortran 2D array.
        \end{itemize}

        The goal is to identify subsets of paralogs whose summed expression vectors approximate $\vec{o}$, within a threshold defined as the lower five percentile of family centroid distances (RDI threshold). Two detection scenarios are supported:

        \begin{enumerate}
            \item \textbf{Subfunctionalization:} The paralogs exhibit significant angles to the ortholog, partitioning ancestral expression.
            \item \textbf{Dosage effect:} The paralogs have small angles to the ortholog and sum to a magnitude significantly exceeding $\|\vec{o}\|$.
        \end{enumerate}

        \subsection{Data Structures}

        \begin{itemize}
            \item All subset membership is encoded as bitmasks of the ranks of the paralog indices, stored as integer values.
            \item A dynamic programming storage array $M$ holds the active subsets to be extended at each iteration, also represented as bitmasks.
            \item An array of norms for each paralog, sorted ascending, is maintained for pruning decisions.
            \item A logical vector marks which paralogs remain candidates at each step.
        \end{itemize}

    \subsection{Algorithm Steps}

        \paragraph{Initialization}
        \begin{itemize}
            \item For all single paralogs, measure their angle to $\vec{o}$.
            \item Apply angle threshold filters:
            \begin{itemize}
                \item Subfunctionalization: only paralogs with angles greater than the gene-family median angle are marked active.
                \item Dosage effect: only paralogs with angles below the gene-family median or lower five percentile are marked active.
            \end{itemize}
            \item Generate all pairs (k=2) of active paralogs to begin.
        \end{itemize}

        \paragraph{Dynamic Programming Iteration}
        For each iteration $k$:
        \begin{enumerate}
            \item Iterate over the bitmask-encoded subsets of size $k-1$ stored in $M$.
            \item For each such subset:
                \begin{enumerate}
                    \item Extend the subset by adding one more paralog, using bitwise operations to set zero bits to one.
                    \item Compute the sum of the expression vectors of the extended subset.
                    \item Compute the residual:
                        \[
                        \vec{r} = \vec{o} - \sum_{i \in S} \vec{p}_i
                        \]
                    \item Test whether:
                        \[
                        \|\vec{r}\| < \text{RDI threshold}
                        \]
                    If true, store this bitmask in the results array and \emph{do not} pass it to the next iteration.
                    \item Otherwise, apply pruning logic:
                    \item For \emph{subfunctionalization}:
                          \begin{itemize}
                              \item Check the minimal length of any still-active paralog:
                                    \[
                                    \min_{i} \|\vec{p}_i\| < \|\vec{r}\|
                                    \]
                                    If this fails, prune the subset as further extension is impossible.
                          \end{itemize}
                    \item For \emph{dosage effect}:
                          \begin{itemize}
                              \item For dosage effect, measure the angle of the current subset sum to $\vec{o}$ and prune if this exceeds the acceptable bound.
                          \end{itemize}
                    \item If the subset is marked to continue, store its bitmask in $M$ for the next iteration.
                \end{enumerate}
            \item Increment $k$ and repeat until $k$ exceeds the maximum paralog search size or no active subsets remain.
        \end{enumerate}

        \paragraph{Final Filtering}
        At the end of the dynamic program, all positive hits are stored as bitmask-encoded subsets. These represent:
        \begin{itemize}
            \item Subfunctionalized paralog subsets whose summed expression reconstructs $\vec{o}$ within the RDI threshold.
            \item Dosage-effect subsets whose sum is aligned with $\vec{o}$ but exceeds its magnitude.
        \end{itemize}

    \subsection{Notes on Implementation}

        \begin{itemize}
            \item Subset bitmasks can be efficiently enumerated by setting zero bits in rank order to one.
            \item Final validation should always be performed using the RDI threshold, regardless of pruning.
            \item A global maximum paralog subset size (e.g., 15) is recommended to control complexity.
            \item For large families, consider pre-filtering based on tissue specificity, expression variance, or percentile angles.
        \end{itemize}

    \subsection{Pseudocode}

        \begin{verbatim}
Initialize:
results = []
active_sets = all bitmask pairs from angle-filtered paralogs

for k = 2 to MAX_PARALOGS:
    next_active_sets = []
    for each bitmask in active_sets:
        for each still-active paralog not in bitmask:
            candidate_mask = bitmask with paralog bit set
            sum_vector = sum of paralogs in candidate_mask
            residual = o - sum_vector
            if norm(residual) < threshold:
                store candidate_mask in results
            else:
                if dosage:
                    if angle(sum_vector, o) too large:
                        continue
                if min(paralog norms) >= norm(residual):
                    continue
                next_active_sets.append(candidate_mask)
    active_sets = next_active_sets
\end{verbatim}

        \noindent
        This algorithm achieves robust and scalable detection of subfunctionalization and dosage patterns by systematic, bitmask-driven dynamic programming with domain-motivated pruning rules.